\section{Ordinary vectors and probability limits}\label{sec:ordinary}

We first control vectors in the lattices themselves.  A count near the
origin will prevent escape into the cusp, and a count near a vector with
no zero coordinate will bound the mass of small neighborhoods in \(X_4\).
The latter bound is the input to the entropy argument in
Section~\ref{sec:rigidity}.  We follow the method of
\cite[Sections~2--5]{PrimeDegreePackets}, proving the estimates with
constants uniform over primitive totally real quartic fields, their orders,
and the orderings of their embeddings.  The absence of a quadratic
subfield supplies the residue estimate needed in this degree.

\subsection{Containment probabilities and unfolding}

For a rational prime \(p\), write
\[
 R_p=\prod_{\mathfrak p\mid p}R_{\mathfrak p}.
\]
Here \(R_{\mathfrak p}\) is the integer ring of the completion
\(K_{\mathfrak p}\).  Let \(v_{\mathfrak p}\) be its integer valuation,
normalized to be \(1\) on a uniformizer, and put
\[
 e_{\mathfrak p}=v_{\mathfrak p}(p),\qquad
 f_{\mathfrak p}=[R/\mathfrak p:\Z/p\Z].
\]
Thus \(\Nm\mathfrak p=p^{f_{\mathfrak p}}\) and
\(\sum_{\mathfrak p\mid p}e_{\mathfrak p}f_{\mathfrak p}=4\);
see \cite[Theorem~3.34]{MilneANT}.  For \(p\mid q\) and a tuple
\(\boldsymbol l_p=(l_{\mathfrak p})_{\mathfrak p\mid p}\) of
nonnegative integers, call the set of \(x\in R_p\) with
\(v_{\mathfrak p}(x_{\mathfrak p})=l_{\mathfrak p}\) in every component
the valuation shell \(\boldsymbol l_p\).  Choose \(x\) in this shell
and define
\begin{equation}\label{eq:ord-local-probability}
 P_p(\boldsymbol l_p)
   =\P_{t_p\in R_p^\times}\{x\in t_p\OO_p\},
\end{equation}
where the probability is normalized unit Haar measure.  The definition
is independent of \(x\): the unit group acts transitively on each set
with specified component valuations.

For a nonzero integral \(R\)-ideal \(\mathfrak a\), set
\begin{equation}\label{eq:ord-ideal-weight}
 P(\mathfrak a)=
   \prod_{p\mid q}
   P_p\bigl((v_{\mathfrak p}(\mathfrak a))_{\mathfrak p\mid p}\bigr).
\end{equation}
An empty product, including the case \(q=1\), is \(1\).  If
\(0\ne b\in I\), its image has no zero local component and
\((b)I^{-1}\) is integral.  In the coordinates
\(I_p=\xi_pR_p\), the valuations of \(\xi_p^{-1}b\) are those of
\((b)I^{-1}\).  Lemma~\ref{lem:pkt-local-labels} and the independence
of the local choices therefore give
\begin{equation}\label{eq:ord-membership}
 \P_{L\in\mathscr L_{\OO}(I)}\{b\in L\}
       =P((b)I^{-1}).
\end{equation}

For a nonnegative Borel function \(f\) on \(\R^4\), write
\[
 E_f(\Lambda)=\sum_{0\ne x\in\Lambda}f(x).
\]
This is a Borel function on \(X_4\): in a local basis chart it is a
countable sum of the Borel functions obtained by evaluating \(f\) at
the integer linear combinations of that basis.
For \(z>0\), define the integral over the surface with absolute
coordinate product \(z\) by
\begin{equation}\label{eq:ord-fiber}
 \mathcal V_f(z)=
 \sum_{\omega\in\{\pm1\}^4}
 \int_{\R^3}
 f\bigl(\omega_1e^{u_1},\omega_2e^{u_2},\omega_3e^{u_3},
        \omega_4z e^{-u_1-u_2-u_3}\bigr)\,du_1\,du_2\,du_3 .
\end{equation}
The signs and logarithmic measure are exactly those of \(T_\infty\).

\begin{proposition}[Ordinary vector unfolding]\label{prop:ord-unfolding}
For every nonnegative Borel \(f\) on \(\R^4\),
\begin{equation}\label{eq:ord-unfolding}
 \int_{X_4}E_f\,d\mu_*
 =\frac1{\sqrt D\,\kappa}
   \sum_{\substack{\mathfrak a\subset R\\\mathfrak a\ne0}}
       P(\mathfrak a)\,
       \mathcal V_f\bigl(\Nm\mathfrak a/X\bigr).
\end{equation}
Both sides are allowed to be infinite.
\end{proposition}

\begin{proof}
This is the unfolding of \cite[Proposition~2.2]{PrimeDegreePackets}
with the packet normalization of Proposition~\ref{prop:pkt-model}.
Average the vector sum over \(L\) in
Equation~\eqref{eq:pkt-integration}, and use
Equation~\eqref{eq:ord-membership}.  The factor \(N_{\OO}\) cancels,
and the left side becomes
\[
 \frac1{h_K\vol(\mathscr F_U)}
 \sum_I\int_{\mathscr F_U}
   \sum_{0\ne b\in I}P((b)I^{-1})f(c_I\sigma(b)v)\,dv .
\]
For fixed \(I\), the map \(b\mapsto(b)I^{-1}\) takes \(I\setminus\{0\}\)
onto the integral ideals in the inverse class of \(I\).  Each fiber
is one free \(U\)-orbit.  Indeed, two generators of the same principal
ideal differ by a unique unit.  The coefficient \(P((b)I^{-1})\) is
constant on that orbit.  Summing its unit translates and integrating
over \(\mathscr F_U\) gives
\[
 \int_{\mathscr F_U}\sum_{\epsilon\in U}
       f(c_I\sigma(\epsilon b)v)\,dv
     =\int_{T_\infty}f(c_I\sigma(b)v)\,dv .
\]
The translates tile \(T_\infty\), including its sign components.
All changes in the order of integration and summation are justified
by Tonelli's theorem.

The absolute product of the coordinates of \(c_I\sigma(b)\) is
\[
 c_I^4|\Nm_{K/\Q}(b)|
 =\frac{\Nm((b)I^{-1})}{X}.
\]
Translation of the first three logarithmic coordinates and relabeling
of the signs therefore identify the last integral with
\(\mathcal V_f(\Nm((b)I^{-1})/X)\), with Jacobian one.  Each integral
ideal belongs to the inverse class of exactly one chosen \(I\).
Equation~\eqref{eq:pkt-unit-volume} now gives the coefficient in
Equation~\eqref{eq:ord-unfolding}.
\end{proof}

\subsection{The mass and tail of the local weights}

The ideal sum in Equation~\eqref{eq:ord-unfolding} need not have simple
coefficients at primes dividing \(q\).  We will use its exact total
mass over valuation shells and a positive moment of their norms.
Both statements are the order case of
\cite[Lemma~3.1]{PrimeDegreePackets}; the argument below explains why
they are uniform even when the local order has large index.

For \(p\mid q\), put
\[
 q_p=[R_p:\OO_p],\qquad
 N_p(\boldsymbol l_p)=p^{\sum_{\mathfrak p\mid p}
                                  f_{\mathfrak p}l_{\mathfrak p}},
 \qquad
 Z_p=\prod_{\mathfrak p\mid p}(1-p^{-f_{\mathfrak p}})^{-1}.
\]
For a collection \(\boldsymbol l=(\boldsymbol l_p)_{p\mid q}\), write
\[
 d(\boldsymbol l)=\prod_{p\mid q}N_p(\boldsymbol l_p),\qquad
 P(\boldsymbol l)=\prod_{p\mid q}P_p(\boldsymbol l_p),\qquad
 Z_q=\prod_{p\mid q}Z_p .
\]
Thus \(d(\boldsymbol l)\) is the norm of the part of an integral ideal
supported above primes dividing \(q\), and \(\prod_{p\mid q}q_p=q\).

\begin{lemma}[Local mass and moment]\label{lem:ord-shells}
With the preceding notation,
\begin{equation}\label{eq:ord-shell-mass}
 \sum_{\boldsymbol l}\frac{P(\boldsymbol l)}{d(\boldsymbol l)}
       =\frac{Z_q}{q}.
\end{equation}
For every \(\eta>0\) there is a constant \(C_\eta\), independent of
\(K,\OO\), such that
\begin{equation}\label{eq:ord-shell-moment}
 \sum_{\boldsymbol l}
       P(\boldsymbol l)d(\boldsymbol l)^{-1+1/8}
       \le C_\eta\,\frac{q^\eta}{q}.
\end{equation}
In particular there is an absolute constant \(C\) for which
\begin{equation}\label{eq:ord-shell-tail}
 \sum_{d(\boldsymbol l)>X^{1/8}}
       \frac{P(\boldsymbol l)}{d(\boldsymbol l)}
       \le \frac Cq X^{-1/128}.
\end{equation}
\end{lemma}

\begin{proof}
Normalize additive Haar measure \(m_p\) on \(R_p\) by \(m_p(R_p)=1\).
In a component \(R_{\mathfrak p}\), the set of elements of valuation
\(l_{\mathfrak p}\) has measure
\((1-p^{-f_{\mathfrak p}})p^{-f_{\mathfrak p}l_{\mathfrak p}}\).
Consequently the shell \(\boldsymbol l_p\) in \(R_p\) has measure
\begin{equation}\label{eq:ord-shell-measure}
 m_p(\text{shell }\boldsymbol l_p)
       =\frac1{Z_p\,N_p(\boldsymbol l_p)}.
\end{equation}
Its conditional additive measure is invariant under the transitive
unit action.  It is therefore the image of unit Haar probability.
By Equation~\eqref{eq:ord-local-probability},
\(P_p(\boldsymbol l_p)\) is exactly the fraction of this shell lying
in \(\OO_p\).  The elements with a zero component have measure zero,
and \(m_p(\OO_p)=q_p^{-1}\).  Summing
Equation~\eqref{eq:ord-shell-measure} over the intersections with
\(\OO_p\) gives
\[
 \sum_{\boldsymbol l_p}
       \frac{P_p(\boldsymbol l_p)}{N_p(\boldsymbol l_p)}
       =\frac{Z_p}{q_p}.
\]
Multiplication over \(p\mid q\) proves
Equation~\eqref{eq:ord-shell-mass}.

For the moment, choose \(x\) with normalized additive Haar probability
on \(\OO_p\), and let \(v_{\mathfrak p}=v_{\mathfrak p}(x_{\mathfrak p})\).
The projection \(H_{\mathfrak p}\) of \(\OO_p\) into
\(R_{\mathfrak p}\) contains the embedded copy of \(\Z_p\), because
\(1\in\OO_p\).  Its induced probability is normalized Haar on
\(H_{\mathfrak p}\).  For an integer \(k\ge1\), the image of this
copy of \(\Z_p\) modulo the ideal of valuation \(k\) has
\(p^{\lceil k/e_{\mathfrak p}\rceil}\) elements.  Hence the image of
\(H_{\mathfrak p}\) has at least that many elements, and
\begin{equation}\label{eq:ord-valuation-tail}
 \P_{\OO_p}\{v_{\mathfrak p}\ge k\}
       \le p^{-\lceil k/e_{\mathfrak p}\rceil}.
\end{equation}
This uses only the size of the additive image of the projection; it
does not require the projections to be independent.

Set
\[
 j=\max_{\mathfrak p\mid p}
              \left\lceil\frac{v_{\mathfrak p}}{e_{\mathfrak p}}\right\rceil.
\]
If \(j\ge a\ge1\), then some \(v_{\mathfrak p}\) is at least
\(e_{\mathfrak p}(a-1)+1\).  Equation~\eqref{eq:ord-valuation-tail}
and a union bound over at most four components give
\(\P_{\OO_p}\{j\ge a\}\le4p^{-a}\).  On the other hand,
\[
 N_p(x)=p^{\sum_{\mathfrak p\mid p}f_{\mathfrak p}v_{\mathfrak p}}
       \le p^{j\sum_{\mathfrak p\mid p}e_{\mathfrak p}f_{\mathfrak p}}
       =p^{4j}.
\]
These assertions hold outside the null sets where a component vanishes.
The tail formula for the expectation of \(p^{j/2}\) now yields
\begin{align}
 \E_{\OO_p}N_p(x)^{1/8}
 &\le \E_{\OO_p}p^{j/2} \notag\\
 &=1+\sum_{a\ge1}
       (p^{a/2}-p^{(a-1)/2})\P_{\OO_p}\{j\ge a\}
 \le 1+4p^{-1/2}.
 \label{eq:ord-local-moment}
\end{align}

Applying Equation~\eqref{eq:ord-shell-measure} once more gives the
exact factorization
\[
 \sum_{\boldsymbol l}
       P(\boldsymbol l)d(\boldsymbol l)^{-1+1/8}
 =\frac1q\prod_{p\mid q}
       \left(Z_p\,\E_{\OO_p}N_p(x)^{1/8}\right).
\]
Each factor is at most
\((1-p^{-1})^{-4}(1+4p^{-1/2})\).  For any fixed \(\eta>0\), this is
at most \(p^\eta\) at all sufficiently large primes.  The product of
the factors at the finitely many remaining primes is bounded by a
constant depending only on \(\eta\).  Since
\(\prod_{p\mid q}p\le q\), this proves
Equation~\eqref{eq:ord-shell-moment}.  Finally, take \(\eta=1/128\)
there and use \(q\le X\):
\[
 \sum_{d(\boldsymbol l)>X^{1/8}}\frac{P(\boldsymbol l)}{d(\boldsymbol l)}
 \le X^{-1/64}
       \sum_{\boldsymbol l}P(\boldsymbol l)d(\boldsymbol l)^{-1+1/8}
 \le \frac Cq X^{-1/128}.
\]
This proves Equation~\eqref{eq:ord-shell-tail}.
\end{proof}

\subsection{Ideal counts relative to the residue}

The shell identity supplies the factor \(q^{-1}\).  To use it in
Equation~\eqref{eq:ord-unfolding}, the count of the remaining ideal
part must retain the factor \(\kappa\).  Recall
\[
 a_K(n)=\#\{\mathfrak a\subset R:\mathfrak a\ne0,\ \Nm\mathfrak a=n\},
 \qquad n\ge1,
\]
and let \(d_4\) be the fourfold divisor function.  Unique factorization
of ideals makes \(a_K\) multiplicative.  At every rational prime,
including a ramified one, its local generating series is
\[
 \sum_{j\ge0}a_K(p^j)T^j
       =\prod_{\mathfrak p\mid p}(1-T^{f_{\mathfrak p}})^{-1}.
\]
Each factor is coefficientwise at most \((1-T)^{-1}\), and there are
at most four factors.  Hence \(a_K(n)\le d_4(n)\) for every \(n\);
in particular \(a_K(p)\le4\).

\begin{lemma}[Residue and prime sums]\label{lem:ord-residue}
There are absolute constants \(D_0>1\), \(h_0>0\), \(Y_0>1\), and
\(C\ge1\) with the following properties for every primitive totally
real quartic field \(K\).
\begin{enumerate}
\item If \(D\ge D_0\), then
\begin{equation}\label{eq:ord-residue-comparison}
 C^{-1}\kappa\le h\zeta_K(1+h)\le C\kappa
       \qquad(0<h\le4/\log D).
\end{equation}
If \(D<D_0\), the same inequalities hold for \(0<h\le h_0\).
\item One has
\begin{equation}\label{eq:ord-residue-lower}
 \kappa^{-1}\le C\log(3D).
\end{equation}
\item Whenever \(Y\ge\max\{Y_0,D^{1/4}\}\),
\begin{equation}\label{eq:ord-prime-sum}
 \exp\left(\sum_{p\le Y}\frac{a_K(p)}p\right)
       \le C\kappa\log Y .
\end{equation}
\end{enumerate}
The sum in Equation~\eqref{eq:ord-prime-sum} is over rational primes.
\end{lemma}

\begin{proof}
We first treat large \(D\), and put \(L=\log D\).
Stark's zero theorem \cite[Lemma~3]{Stark1974} gives at most one
nontrivial zero of \(\zeta_K\) in
\[
 \Re s\ge1-\frac1{4L},\qquad |\Im s|\le\frac1{4L};
\]
any such zero is real and simple.  His
Lemma~8 \cite[Lemma~8]{Stark1974}, valid also for nonnormal fields,
states in degree four that a real zero
\(\beta\ge1-(96L)^{-1}\) would force a quadratic subfield of \(K\).
Primitivity excludes that subfield.  It follows that every nontrivial
zero \(\rho\) satisfies
\begin{equation}\label{eq:ord-zero-disk}
 |1-\rho|\ge \frac{c_0}{L},\qquad c_0=\frac1{192}.
\end{equation}
Indeed, a zero in the smaller open disk would lie in Stark's rectangle,
would be real, and would be in the forbidden real interval.

For real \(s>1\), the logarithmic derivative of the completed zeta
function gives
\begin{equation}\label{eq:ord-zero-sum}
\begin{aligned}
 S_K(s):=\sum_\rho\Re\frac1{s-\rho}
 &=\frac1s+\frac1{s-1}+\frac L2+
       \frac{\zeta_K'}{\zeta_K}(s)+G(s),\\
 G(s)&=2\left(\frac{\Gamma'(s/2)}{\Gamma(s/2)}-\log\pi\right);
\end{aligned}
\end{equation}
see \cite[Equation~(9)]{Stark1974}.  The real zero sum converges and
has nonnegative terms, since \(\Re\rho\le1\).  At \(s_0=1+L^{-1}\),
the Euler product gives \(\zeta_K'/\zeta_K(s_0)\le0\).  The function
\(G\) is bounded on \([1,2]\), so for sufficiently large \(D\),
\[
 0\le S_K(s_0)\ll L.
\]
This estimate persists for all \(1<s\le1+4/L\).  To see this term by
term, write \(\rho=\beta+i\gamma\).  Since \(\beta\le1\),
\[
 s-\beta\le4(s_0-\beta),\qquad
 |s-\rho|\ge|1-\rho|.
\]
Equation~\eqref{eq:ord-zero-disk} consequently gives
\[
 |s_0-\rho|\le |s-\rho|+\frac3L
              \le (1+3/c_0)|s-\rho|.
\]
Using \(\Re(s-\rho)^{-1}=(s-\beta)/|s-\rho|^2\), we obtain
\[
 \Re\frac1{s-\rho}
 \le4(1+3/c_0)^2\Re\frac1{s_0-\rho}.
\]
Thus \(S_K(s)\ll L\) on the entire interval.

Let \(\mathcal Z_K(s)=(s-1)\zeta_K(s)\).  Rearranging
Equation~\eqref{eq:ord-zero-sum} shows that
\[
 \frac{\mathcal Z_K'}{\mathcal Z_K}(s)
       =S_K(s)-\frac1s-\frac L2-G(s)=O(L)
       \qquad(1<s\le1+4/L).
\]
The function \(\mathcal Z_K\) extends to \(s=1\) with value
\(\kappa>0\).  Integrating its logarithmic derivative across an
interval of length at most \(4/L\) proves
Equation~\eqref{eq:ord-residue-comparison}, with an absolute constant.
At \(s_0\), the inequality \(\zeta_K(s_0)\ge1\) then gives
\(\kappa\gg L^{-1}\).

There are only finitely many fields with \(D<D_0\), by Hermite's
finiteness theorem \cite[Theorem~8.43]{MilneANT}.  For each of them,
\(h\zeta_K(1+h)\) tends to the positive residue \(\kappa\) as
\(h\downarrow0\).  Taking a common sufficiently small \(h_0\) over
this finite set proves the remaining comparison, and their positive
residues have a common positive lower bound.  Enlarging \(C\) proves
Equation~\eqref{eq:ord-residue-lower} for every field.

Choose \(Y_0\) sufficiently large that \(1/\log Y_0\le h_0\), and
let \(Y\ge\max\{Y_0,D^{1/4}\}\).  For \(h=1/\log Y\), the comparison
just proved applies: in the large \(D\) case \(h\le4/\log D\), and
in the bounded \(D\) case \(h\le h_0\).  Positivity in the Euler
logarithm gives
\[
 \sum_{p\le Y}a_K(p)p^{-1-h}\le\log\zeta_K(1+h).
\]
Moreover, \(a_K(p)\le4\), \(1-p^{-h}\le h\log p\), and the elementary
Chebyshev bound followed by partial summation imply
\[
 \sum_{p\le Y}a_K(p)(p^{-1}-p^{-1-h})
 \le4h\sum_{p\le Y}\frac{\log p}{p}
 \ll h\log Y=1.
\]
It follows that
\[
 \sum_{p\le Y}\frac{a_K(p)}p
 \le \log\zeta_K(1+h)+O(1)
 \le \log(\kappa\log Y)+O(1),
\]
which proves Equation~\eqref{eq:ord-prime-sum}.  There is no upper
restriction on \(Y\).  This observation is what permits \(D\) to stay
bounded while the order index, and hence \(X\), grows.
\end{proof}

We also need an estimate for ideals in an interval.  The following
uniform form of Shiu's bound \cite{Shiu1980} is the one recorded in
\cite[Equation~(4.2)]{PrimeDegreePackets}.

\begin{lemma}\label{lem:ord-shiu}
There are absolute constants \(C,u_0\) such that, for every
nonnegative multiplicative function \(b\) with \(b(n)\le d_4(n)\),
every real \(u\ge u_0\), and every \(u^{1/2}\le v\le u\),
\begin{equation}\label{eq:ord-shiu}
 \sum_{u-v<n\le u}b(n)
 \le C\frac v{\log u}
       \exp\left(\sum_{p\le u}\frac{b(p)}p\right).
\end{equation}
\end{lemma}

\begin{proof}
Apply \cite[Theorem~1.1]{Pollack2020} with no excluded residue classes
and with its parameter \(\beta=1/4\).  Its two growth conditions hold
with bounds independent of \(b\):
\[
 b(p^j)\le d_4(p^j)=\binom{j+3}{3}\le4^j,\qquad
 b(n)\le d_4(n)\ll_\epsilon n^\epsilon
 \quad(\epsilon>0).
\]
The theorem's constant and threshold depend only on these growth
bounds and on the fixed parameter.  Its strict length condition
\(u^\beta<v\) is satisfied even at \(v=u^{1/2}\) once \(u>1\).
This gives Equation~\eqref{eq:ord-shiu}, including both endpoints.
\end{proof}

\begin{proposition}[Weighted ideal intervals]\label{prop:ord-weighted-intervals}
There is an absolute constant \(C\) such that, for each fixed
\(A\ge1\) and \(c>0\), there is a threshold \(X_0(A,c)\) with the
following property.  For every \(K,\OO,\sigma\) under consideration
with \(X\ge X_0(A,c)\), and all real \(x,y\) satisfying
\[
 X^{3/4}\le x\le AX,\qquad cx\le y\le x,
\]
one has
\begin{equation}\label{eq:ord-weighted-interval}
 \sum_{\substack{\mathfrak a\subset R,\ \mathfrak a\ne0\\
                  x-y<\Nm\mathfrak a\le x}}P(\mathfrak a)
       \le C\,\frac{\kappa y}{q}.
\end{equation}
The constant \(C\) is independent of \(A,c\); the threshold may
depend on them.
\end{proposition}

\begin{proof}
The range is empty for \(c>1\), so suppose \(0<c\le1\).
We adapt the separation of the local norm from the remaining ideal
in \cite[Proposition~4.2]{PrimeDegreePackets}.  Put
\[
 b(n)=a_K(n)\mathbf 1_{(n,q)=1}.
\]
This is nonnegative and multiplicative, and it is at most \(d_4\).
Every integral ideal splits uniquely into its part above the primes
dividing \(q\), described by one tuple \(\boldsymbol l\), and a part
coprime to \(qR\).  The left side of
Equation~\eqref{eq:ord-weighted-interval} is therefore
\begin{equation}\label{eq:ord-ideal-split}
 \sum_{\boldsymbol l}P(\boldsymbol l)
       \sum_{(x-y)/d(\boldsymbol l)<n\le x/d(\boldsymbol l)}b(n).
\end{equation}

First take \(d=d(\boldsymbol l)\le X^{1/8}\), and put \(u=x/d\),
\(v=y/d\).  Then
\[
 u\ge X^{5/8}\ge D^{1/4},\qquad v\ge cu.
\]
The second inequality implies \(v\ge u^{1/2}\) when \(X\) is large
in terms of \(c\).  The first inequality puts \(u\) in the range of
Equation~\eqref{eq:ord-prime-sum}, once \(X\) is above an absolute
threshold.  To remove the primes dividing \(q\) from that Euler bound,
observe that
\begin{equation}\label{eq:ord-euler-removal}
 \log Z_q=\sum_{p\mid q}\frac{a_K(p)}p+O(1),\qquad
 \sum_{\substack{p\mid q\\p>u}}\frac{a_K(p)}p
       \ll\frac{\log(2q)}u\ll1 .
\end{equation}
For the first identity, expand
\(\log Z_p=\sum_{\mathfrak p\mid p}\sum_{j\ge1}p^{-jf_{\mathfrak p}}/j\).
The terms of order \(p^{-1}\) count exactly the primes of residue
degree one, and the remaining terms are \(O(p^{-2})\), uniformly
in the splitting type.  Their sum over rational primes is bounded.
For the second estimate, use \(a_K(p)\le4\), at most
\(\log(2q)/\log2\) primes dividing \(q\), and
\(u\ge X^{5/8}\), \(q\le X\).
Consequently Equations~\eqref{eq:ord-prime-sum} and
\eqref{eq:ord-euler-removal} give
\[
 \exp\left(\sum_{p\le u}\frac{b(p)}p\right)
       \ll \frac{\kappa\log u}{Z_q}.
\]
Lemma~\ref{lem:ord-shiu} and
Equation~\eqref{eq:ord-shell-mass} now bound the contribution of
these tuples to Equation~\eqref{eq:ord-ideal-split} by
\[
 C\frac{\kappa y}{Z_q}
       \sum_{d(\boldsymbol l)\le X^{1/8}}
             \frac{P(\boldsymbol l)}{d(\boldsymbol l)}
 \le C\frac{\kappa y}{q}.
\]

It remains to control tuples with \(d>X^{1/8}\).  Those with \(d>x\)
contribute nothing.  For \(z\ge1\), summing the first three factors
of the divisor convolution gives
\[
 \sum_{n\le z}d_4(n)
 =\sum_{a_1a_2a_3a_4\le z}1
 \le z\left(\sum_{a\le z}\frac1a\right)^3
 \le z(1+\log z)^3 .
\]
Use this with \(z=x/d\), and then
Equation~\eqref{eq:ord-shell-tail}.  The remaining contribution is
at most
\[
 Cx(1+\log x)^3
       \sum_{d(\boldsymbol l)>X^{1/8}}
             \frac{P(\boldsymbol l)}{d(\boldsymbol l)}
 \le C\frac{x}{q}X^{-1/128}(1+\log x)^3 .
\]
For fixed \(A\), once \(X\ge A\) the condition \(x\le AX\) gives
\(\log x\le2\log X\).  By
Equation~\eqref{eq:ord-residue-lower} and \(D\le X^2\), the ratio of
this last bound to \(\kappa x/q\) is at most
\[
 C X^{-1/128}(\log(3X))^4=o(1).
\]
This bound is independent of the field and the order.  Increase the
threshold in terms of \(A,c\) until the ratio is at most \(c\).
The tail is then at most \(\kappa cx/q\le\kappa y/q\).  The constant
in the resulting bound has not changed with \(A\) or \(c\), proving
the proposition.
\end{proof}

\subsection{Short vectors and small neighborhoods}

We now apply the weighted interval bound to the two vector counts
announced at the start of the section.  If \(f_r=\mathbf 1_{[-r,r]^4}\), direct integration
in Equation~\eqref{eq:ord-fiber} gives
\begin{equation}\label{eq:ord-cube-fiber}
 \mathcal V_{f_r}(z)=\frac{16}{6}
                       \bigl(\log_+(r^4/z)\bigr)^3 .
\end{equation}
Indeed, in a fixed sign component put \(s_i=\log r-u_i\) for
\(1\le i\le3\).  The cube conditions become \(s_i\ge0\) and
\(\sum_{i=1}^3s_i\le\log(r^4/z)\), a simplex of the displayed volume
divided by sixteen.

\begin{proposition}[Tightness]\label{prop:ord-tightness}
There is an absolute constant \(C\) such that, for every fixed
\(0<r<1\),
\begin{equation}\label{eq:ord-cusp}
 \limsup_{X\longrightarrow\infty}
 \mu_*\{\Lambda:\text{some }0\ne v\in\Lambda
                         \text{ satisfies }\|v\|_\infty<r\}
       \le Cr^4 .
\end{equation}
The limit superior ranges over all primitive totally real quartic
fields, all their orders, and all orderings of their real embeddings.
Consequently, for every sequence of these packet measures
\(\mu_i=\mu_{\OO_i,\sigma_i}\) with \(X_i\to\infty\) and every
\(\epsilon>0\), there is a compact \(\mathcal K_\epsilon\subset X_4\) with
\(\mu_i(\mathcal K_\epsilon)\ge1-\epsilon\) for all sufficiently large \(i\).
Such a sequence is tight, and every weak subsequential limit is an
\(A_4\)-invariant probability on \(X_4\).
\end{proposition}

\begin{proof}
The count \(E_{f_r}\) dominates the indicator in
Equation~\eqref{eq:ord-cusp}.  Put \(T=r^4X\) and \(T_0=X^{3/4}\).
By Equations~\eqref{eq:ord-unfolding} and \eqref{eq:ord-cube-fiber},
it suffices to bound
\begin{equation}\label{eq:ord-cube-sum}
 \frac1{\sqrt D\,\kappa}
 \sum_{\Nm\mathfrak a\le T}
       P(\mathfrak a)
       \bigl(\log(T/\Nm\mathfrak a)\bigr)^3 .
\end{equation}
The sum is empty when \(T<1\).  Otherwise its logarithmic factor is
at most \((\log X)^3\), since \(\Nm\mathfrak a\ge1\) and \(r<1\).
Apply Proposition~\ref{prop:ord-weighted-intervals} at
\(x=y=T_0\), with \(A=c=1\).  The part with
\(\Nm\mathfrak a\le\min\{T,T_0\}\) is at most
\[
 C\frac{\kappa T_0}{q}(\log X)^3.
\]
After division by \(\sqrt D\,\kappa=X\kappa/q\), this contributes
\(O(X^{-1/4}(\log X)^3)=o(1)\).

If \(T>T_0\), put
\[
 A_{T_0}(t)=\sum_{T_0<\Nm\mathfrak a\le t}P(\mathfrak a)
 \qquad(T_0\le t\le T).
\]
The same interval estimate at \(x=y=t\) gives
\(A_{T_0}(t)\le C\kappa t/q\), since \(T\le X\).  Writing each
logarithmic cube as an integral and using Tonelli yields
\begin{align*}
 \sum_{T_0<\Nm\mathfrak a\le T}
       P(\mathfrak a)\bigl(\log(T/\Nm\mathfrak a)\bigr)^3
 &=3\int_{T_0}^T A_{T_0}(t)
                   \bigl(\log(T/t)\bigr)^2\frac{dt}{t}\\
 &\le C\frac{\kappa}{q}
              \int_0^T\bigl(\log(T/t)\bigr)^2\,dt
 \le C\frac{\kappa T}{q}.
\end{align*}
The last integral is finite after the substitution \(t=Ts\).
This part of Equation~\eqref{eq:ord-cube-sum} is therefore \(O(r^4)\).
If \(T\le T_0\), it is absent.  The constants used at \(A=c=1\) are
absolute, so the resulting estimate proves
Equation~\eqref{eq:ord-cusp}.  This is the compactness argument of
\cite[Proposition~5.1]{PrimeDegreePackets}, now using the quartic
weighted interval bound.

Mahler's compactness criterion \cite[Section~3, Theorem~2]{Mahler1946}
says that
\[
 \mathcal K_r=\{\Lambda\in X_4:\|v\|_\infty\ge r
                                      \text{ for all }0\ne v\in\Lambda\}
\]
is compact.  Choose \(r\) so that \(Cr^4<\epsilon/2\).
Equation~\eqref{eq:ord-cusp} then gives
\(\mu_i(\mathcal K_r)\ge1-\epsilon\) for all sufficiently large \(i\).
Each of the finitely many preceding packet measures is supported on a
finite union of compact orbits, by Proposition~\ref{prop:pkt-model}.
Enlarging the compact set to include those supports proves tightness
of the entire sequence.  The space \(X_4\) is locally compact,
second countable, and metrizable; the weak compactness theorem for
tight probabilities therefore gives probability subsequential limits.
For \(a\in A_4\) and \(f\in C_c(X_4)\), the equality
\[
 \int f(xa)\,d\mu_i(x)=\int f(x)\,d\mu_i(x)
\]
passes to a weak limit, since both test functions are continuous with
compact support.  This proves \(A_4\)-invariance.
\end{proof}

\begin{lemma}[Boxes away from the coordinate hyperplanes]\label{lem:ord-boxes}
Let \(V_0\subset(\R^\times)^4\) be compact, and let \(b>0\).
There are constants \(C,r_0>0\), depending only on \(V_0,b\), such
that for every fixed \(0<r<r_0\) there is a threshold
\(X_0(V_0,b,r)\) for which
\begin{equation}\label{eq:ord-vector-box}
 \int_{X_4}E_{\mathbf 1_{v+[-br,br]^4}}\,d\mu_*
       \le Cr^4\qquad(v\in V_0,\ X\ge X_0(V_0,b,r)).
\end{equation}
The threshold and constants are uniform over \(K,\OO,\sigma\) and
over the centers \(v\in V_0\).
\end{lemma}

\begin{proof}
The assertion is immediate if \(V_0\) is empty.  Otherwise choose
\(0<\epsilon\le M\) such that
\(\epsilon\le|v_i|\le M\) for every \(v\in V_0\) and \(i\).
For \(r_0<\epsilon/(2b)\), each box in the statement lies in one
coordinate sign component.  Its first three logarithmic widths satisfy
\[
 \log\frac{|v_i|+br}{|v_i|-br}\ll_{V_0,b}r .
\]
The product map \(x\mapsto\prod_i|x_i|\) has uniformly bounded first
derivatives on these boxes.  If \(t_v=\prod_i|v_i|\), it follows
from Equation~\eqref{eq:ord-fiber} that, for one \(C_1\ge1\),
\begin{equation}\label{eq:ord-box-fiber}
 \mathcal V_{\mathbf 1_{v+[-br,br]^4}}(z)
       \le C_1r^3\,\mathbf 1_{[t_v-C_1r,t_v+C_1r]}(z)
       \qquad(v\in V_0,\ 0<r<r_0).
\end{equation}
Let \(t_->0\) and \(t_+<\infty\) be the minimum and maximum of \(t_v\)
on \(V_0\), and shrink \(r_0\) so that \(6C_1r_0<t_-\).
For a fixed \(r\) in this range set
\[
 x=(t_v+2C_1r)X,\qquad y=4C_1rX.
\]
The possible ideal norms in Equation~\eqref{eq:ord-box-fiber} lie
strictly inside \((x-y,x)\).  The enlargement avoids any issue with
integer endpoints.  There is a fixed \(A\ge1\), depending only on
\(V_0,b\), such that \(x\le AX\) for all centers and radii under
consideration.  For sufficiently large \(X\), uniformly in the
centers, \(x\ge X^{3/4}\).  Also \(y\le x\) by the choice of \(r_0\),
and \(y\ge cx\) with \(c=4C_1r/A>0\), after decreasing \(r_0\) if
necessary so that \(c\le1\).

Proposition~\ref{prop:ord-weighted-intervals} applies for every fixed
\(r\), with this \(c\) and the fixed \(A\).  Its constant is independent
of \(c\), so the sum of the relevant ideal weights is at most
\(C\kappa rX/q\), with a constant independent of \(r\).  Combining
this with Equations~\eqref{eq:ord-unfolding} and
\eqref{eq:ord-box-fiber}, and using \(\sqrt D=X/q\), proves
Equation~\eqref{eq:ord-vector-box}.  Only the threshold for \(X\)
depends on the fixed radius.
\end{proof}

For \(r>0\), define the open identity neighborhood
\begin{equation}\label{eq:ord-ball-neighborhood}
 B(r)=\{h\in\mathrm{SL}_4(\R):
              \max_{i,j}|h_{ij}-\delta_{ij}|<r\}.
\end{equation}

\begin{proposition}[Local bounds in probability limits]\label{prop:ord-balls}
Let \(\mu\) be a weak probability limit of packet measures along
\(X_i\to\infty\).  For every compact \(\Omega\subset X_4\), there
are constants \(C_\Omega,r_\Omega>0\) such that
\begin{equation}\label{eq:ord-balls}
 \mu(xB(r))\le C_\Omega r^4
       \qquad(x\in\Omega,\ 0<r<r_\Omega).
\end{equation}
\end{proposition}

\begin{proof}
The assertion is immediate if \(\Omega\) is empty, so assume otherwise.
We use the passage from vector boxes to ordinary neighborhoods in
\cite[Proposition~5.2]{PrimeDegreePackets}.  First, every full lattice
has a vector with no zero coordinate.  If \(\Lambda=\Z^4g\), each
coordinate of \((1,N,N^2,N^3)g\) is a nonzero polynomial in \(N\):
its coefficients are a column of the invertible matrix \(g\).
An integer \(N\) outside their finitely many roots gives the required
vector.

For each \(x_0\in\Omega\), choose a continuous local lift \(g(x)\) of
nearby lattices to \(\mathrm{SL}_4(\R)\), and an integer row \(z\)
for which \(zg(x_0)\) has no zero coordinate.  On a sufficiently
small neighborhood with compact closure in this chart, \(zg(x)\)
stays bounded and its coordinates stay bounded away from zero.
A finite collection of these neighborhoods covers \(\Omega\).
Consequently we may select, for every \(x\in\Omega\), a vector
\(v_x\in\Lambda_x\) so that all selected vectors lie in one compact
set \(V_0\subset(\R^\times)^4\).  The selection is used pointwise and
need not be measurable.

Choose \(M\) with \(|(v_x)_i|\le M\) for all selected vectors, and
fix \(b>4M\).  If \(y=xh\) with \(h\in B(r)\), then
\(v_xh\in\Lambda_y\), and
\[
 |(v_xh-v_x)_j|
 \le \sum_i |(v_x)_i|\,|h_{ij}-\delta_{ij}|<br.
\]
Thus every lattice in \(xB(r)\) has a nonzero vector in
\(v_x+[-br,br]^4\).  Lemma~\ref{lem:ord-boxes} gives, for each fixed
small \(r>0\) and all sufficiently large \(i\),
\[
 \mu_i(xB(r))
 \le \int E_{\mathbf 1_{v_x+[-br,br]^4}}\,d\mu_i
 \le C_\Omega r^4
 \qquad(x\in\Omega).
\]
The constants are common to all \(x\) and \(r\) in the stated ranges;
the index after which this estimate holds may depend on \(r\).
The set \(xB(r)\) is open in \(X_4\).  For each fixed \(x,r\),
Portmanteau's inequality therefore yields
\[
 \mu(xB(r))\le\liminf_{i\to\infty}\mu_i(xB(r))
       \le C_\Omega r^4.
\]
This proves the bound for every \(r\) with the same constant.  It uses
no interchange between a shrinking radius and the packet limit.
\end{proof}
